% Generated draft: requires mathematical and source review.
\documentclass[11pt]{article}
\usepackage[margin=25mm]{geometry}
\usepackage{fontspec}
\setmainfont{Latin Modern Roman}
\usepackage{amsmath,amssymb,mathtools,unicode-math}
\setmathfont{Latin Modern Math}
\usepackage{enumitem,needspace,xurl,hyperref}
\setlength{\parindent}{0pt}
\setlength{\parskip}{5pt}
\setlength{\emergencystretch}{3em}
\newcommand{\N}{\mathbb{N}}
\newcommand{\Z}{\mathbb{Z}}
\newcommand{\Q}{\mathbb{Q}}
\newcommand{\R}{\mathbb{R}}
\newcommand{\C}{\mathbb{C}}
\newcommand{\F}{\mathbb{F}}
\newcommand{\A}{\mathbb{A}}
\newcommand{\PP}{\mathbb{P}}
\newcommand{\E}{\mathbb{E}}
\newcommand{\eps}{\varepsilon}
\newcommand{\dd}{\,\mathrm{d}}
\begin{document}
\section*{UnsolvedMath ID: 700662}
\textbf{Source problem code:} catalogue entry
\subsection{Definitions and mathematical statement}
Let \(R\) be an associative commutative ring with identity such that \(2\cdot 1_R\) and \(3\cdot 1_R\) are units. The identity is implicit in the source's requirement that \(2\) and \(3\) have inverses.

The supplied record does not define \(PE_2(R)\). Under the conventional interpretation as the projective elementary group, define
\[
E_2(R)=\left\langle
\begin{pmatrix}1&r\\0&1\end{pmatrix},
\begin{pmatrix}1&0\\r&1\end{pmatrix}
\;\middle|\;r\in R
\right\rangle\leq \operatorname{GL}_2(R),
\]
and
\[
PE_2(R)=E_2(R)\big/\left(E_2(R)\cap\{uI_2:u\in R^\times\}\right),
\]
where \(R^\times\) is the group of units of \(R\) and \(I_2\) is the identity matrix. This interpretation of the notation requires confirmation against the original source.

\textbf{Problem (V. M. Petechuk).} For every such ring \(R\), describe all abstract group automorphisms of \(PE_2(R)\); equivalently, classify all bijections
\[
\varphi:PE_2(R)\longrightarrow PE_2(R)
\quad\text{such that}\quad
\varphi(xy)=\varphi(x)\varphi(y)
\quad\text{for all }x,y\in PE_2(R).
\]
The source requests a description of these automorphisms but supplies no proposed classification or required form of that description. The displayed formulation is conditional on the interpretation of \(PE_2(R)\) above.
\subsection{Short English statement}
For every associative commutative ring in which \(2\) and \(3\) are invertible, describe all automorphisms of its projective elementary group of degree \(2\). The meaning of \(PE_2(R)\) must be confirmed because the supplied record does not define it. The registry labels the problem open; this is supplied metadata, not an independently verified assessment.
\subsection{Sources}
\begin{itemize}
\item Supplied local UnsolvedMath registry record: permanent ID \(700662\), exact source problem code \texttt{catalogue entry}, title \emph{Kourovka Notebook Problem 10.47}. Upstream revision: \texttt{d371ee917ac9647fbb06cd6fc529deec6f3dda82}.
\item Supplied source attribution: \emph{The Kourovka Notebook (Unsolved Problems in Group Theory)}, 10th issue (1986), Problem 10.47; proposed by V. M. Petechuk in 1986.
\item The record also identifies the 21st issue edition (2026), \texttt{arXiv:1401.0300v48}. Its collection description names E. I. Khukhro and V. D. Mazurov as editors of Issues 1--21 (1965--2026). These bibliographic details are supplied attribution and have not been independently checked.
\item The supplied registry status is \texttt{open}. No separate literature assessment, literature-verification date, or external source URL is supplied.
\end{itemize}
\end{document}
